% ============================================================
% sec/01introduccion.tex
%
% OBLIGATORIO: contexto y relevancia, pregunta de investigación
% explícita, hipótesis de trabajo, objetivos con contribución
% esperada, y estructura del documento.
% ============================================================
\section{Introducción}
\label{sec:introduccion}

\subsection{Contexto y relevancia}
\label{subsec:contexto}

La estimación de la probabilidad de victoria a partir del estado
intermedio de una partida constituye una función de valor canónica en la
analítica de juegos competitivos. Brill \textit{et al.} señalan que dicha
cantidad no es una estadística contable ni una magnitud observable, sino
que queda definida por un modelo estimado a partir de datos
\cite{brill_winprob_2026}. Esta propiedad distingue la tarea de otros
problemas de predicción: el desenlace observado de una partida es binario,
mientras que la magnitud que se desea estimar es una probabilidad que
nunca se observa directamente.

\razon{La primera oración fija la naturaleza de la tarea. Si no se declara
que la probabilidad verdadera es inobservable, el lector puede suponer que
existe una etiqueta contra la cual verificar cada predicción, y toda la
discusión posterior sobre calibración pierde sentido.}

En el ámbito de los deportes electrónicos, la utilidad de esta estimación
excede el interés metodológico. Hodge \textit{et al.} observan que, a
diferencia de los deportes tradicionales, numerosos juegos competitivos
carecen de un marcador legible, y que su complejidad dificulta el
seguimiento por parte de audiencias no expertas; una estadística única que
resuma el estado de la partida amplía el atractivo del juego y lo vuelve
más accesible \cite{hodge_winprediction_2021}. Xenopoulos \textit{et al.}
desarrollan la misma línea al estimar la probabilidad de victoria por ronda
y analizar cómo varían los efectos de las características del estado entre
poblaciones de jugadores con distinto nivel de habilidad
\cite{xenopoulos_winprob_2022}.

El Pok\'emon Trading Card Game presenta las condiciones que hacen relevante
este tipo de estimación. \pc{2 o 3 oraciones con datos concretos del
dominio: tamaño de la comunidad competitiva, cantidad de cartas
disponibles, complejidad combinatoria del mazo. Este contenido pertenece al
Eje A y lo aporta Daniel.} El estado de una partida involucra premios
restantes, puntos de vida, energías acumuladas, cartas en mano, evoluciones
y estados alterados, distribuidos entre dos jugadores con información
parcial sobre el mazo y la mano del oponente. Un observador no experto no
puede determinar de un vistazo cuál de los dos jugadores se encuentra en
ventaja.

\subsection{Planteamiento del problema}
\label{subsec:problema}

El problema que este trabajo aborda es la ausencia de un estimador validado
de probabilidad de victoria a partir del estado de juego en el Pok\'emon
Trading Card Game, junto con la ausencia de una caracterización de cuánta
información sobre el desenlace contiene efectivamente dicho estado.

Esta ausencia no es únicamente una carencia de resultados: es también una
carencia metodológica. Los conjuntos de datos derivados de repeticiones de
partidas presentan una estructura de dependencia que invalida los
procedimientos de partición aleatoria habituales. Brill \textit{et al.}
demuestran que, cuando múltiples observaciones comparten el mismo sorteo
del desenlace, el sesgo y la varianza de los estimadores se incrementan y
el tamaño de muestra efectivo se reduce sustancialmente
\cite{brill_winprob_2026}. Figueiredo y Mendes cuantifican el efecto
complementario sobre las métricas reportadas: en conjuntos de imágenes con
alta correlación espaciotemporal, la partición aleatoria infla el
desempeño medido respecto de una partición agrupada de referencia
\cite{figueiredo_leakage_2024}.

\razon{Este párrafo debe existir en la introducción y no solo en la
metodología. El riesgo metodológico propio del Track A es la fuga de
información, y presentarlo desde el planteamiento del problema convierte
el control de fuga en parte de la contribución en lugar de en una
precaución técnica mencionada al final.}

\subsection{Pregunta de investigación}
\label{subsec:pregunta}

\begin{quote}
\textbf{PI:} ¿Cuánta información sobre el desenlace de una partida de
Pok\'emon Trading Card Game contiene un estado de juego aislado; cómo varía
esa capacidad predictiva conforme progresa la partida; y en qué medida se
conserva cuando el modelo se evalúa sobre partidas, mazos y periodos
temporales no observados durante el entrenamiento?
\end{quote}

\noindent\textbf{Condiciones de refutación.} La pregunta admite respuesta
negativa en al menos tres formas: que el modelo no supere a una heurística
de dominio sobre partidas no vistas; que la capacidad predictiva no
presente variación apreciable a lo largo de la partida; o que el desempeño
se degrade hasta niveles no informativos bajo los protocolos de
generalización temporal o por mazo.

\razon{El enunciado exige que la pregunta pueda responderse con evidencia y
admita refutación. Enumerar las condiciones concretas bajo las cuales la
respuesta sería negativa es lo que distingue una pregunta de investigación
de una expectativa.}

\noindent\textbf{Condiciones de medición.} La respuesta se sustenta sobre
métricas calculadas en conjuntos de prueba disjuntos a nivel de episodio,
reportando de forma separada la capacidad de discriminación y la calidad de
la calibración, con intervalos de confianza construidos sobre episodios y
no sobre observaciones individuales.

\subsection{Hipótesis de trabajo}
\label{subsec:hipotesis}

\begin{quote}
\textbf{H1.} El estado de juego contiene información predictiva sobre el
desenlace desde etapas tempranas de la partida, y dicha información se
incrementa de forma monótona conforme progresa la partida.
\end{quote}

\begin{quote}
\textbf{H2.} La capacidad de discriminación se conserva en mayor medida
que la calidad de la calibración cuando el modelo se evalúa sobre
periodos posteriores a los de entrenamiento.
\end{quote}

\begin{quote}
\textbf{H3.} \pc{Hipótesis sobre generalización a mazos no vistos. Su
formulación depende del análisis de concentración de mazos que
corresponde al Eje C.}
\end{quote}

\noindent\textbf{Fundamento de H2.} Xenopoulos \textit{et al.} reportan que
un modelo evaluado fuera del entorno en que fue entrenado presenta un
incremento moderado en la pérdida logarítmica acompañado de una degradación
considerablemente mayor en el error de calibración esperado
\cite{xenopoulos_winprob_2022}. Silva Filho \textit{et al.} ofrecen la
explicación formal: las reglas de puntuación propias se descomponen en un
componente de calibración y uno de refinamiento, y las métricas basadas
exclusivamente en el ordenamiento son invariantes ante transformaciones
monótonas de las probabilidades predichas \cite{silvafilho_calibration_2023}.

\noindent\textbf{Advertencia sobre H1.} Brill \textit{et al.} documentan que
el sesgo del estimador es considerablemente mayor al comienzo de la partida,
por dos razones estructurales: los estados tempranos con desequilibrios
marcados son infrecuentes, y resulta más difícil vincular un estado temprano
con un desenlace determinado al final \cite{brill_winprob_2026}. En
consecuencia, un desempeño reducido en la ventana temprana admite dos
interpretaciones que este diseño debe procurar distinguir: ausencia genuina
de señal, o sesgo del estimador por escasez de observaciones en esa región
del espacio de estados.

\razon{Declarar esta ambigüedad antes de obtener resultados es obligatorio.
Sin ella, un desempeño bajo en turnos tempranos se interpretaría como
refutación de H1 cuando podría deberse al sesgo que Brill documenta.}

\decidir{Si el equipo desea separar ambas interpretaciones, hay que definir
en la Etapa 2 un procedimiento adicional: reportar el conteo de episodios y
de observaciones por ventana de turno, o restringir el análisis temprano a
regiones del espacio de estados con soporte suficiente. La decisión afecta
el diseño experimental y conviene tomarla antes de ejecutar.}

\subsection{Objetivos}
\label{subsec:objetivos}

\noindent\textbf{Objetivo general.} Determinar y caracterizar la capacidad
predictiva del estado de juego sobre el desenlace de una partida de
Pok\'emon Trading Card Game, mediante modelos de aprendizaje supervisado
sobre datos tabulares, bajo un protocolo experimental que controle de forma
explícita la fuga de información derivada de la dependencia entre
observaciones.

\noindent\textbf{Objetivos específicos.}
\begin{enumerate}
  \item Formalizar la estimación de la probabilidad de victoria a partir
        del estado de juego como un problema de clasificación binaria
        probabilística sobre datos tabulares, estableciendo la distinción
        entre la unidad de observación y la unidad estadísticamente
        independiente.
  \item Construir un conjunto de datos tabular a partir de las
        repeticiones disponibles, definiendo un diccionario de
        características restringido a la información accesible en el
        instante de la decisión.
  \item Cuantificar el riesgo de fuga de información asociado a la
        partición aleatoria sobre este conjunto de datos, y establecer un
        protocolo de partición agrupada verificable de forma automática.
  \item Implementar y comparar una escalera de modelos de referencia que
        incluya un predictor trivial, una heurística de dominio y modelos
        tabulares clásicos, evaluando de forma separada discriminación y
        calibración.
  \item Caracterizar la evolución de la capacidad predictiva a lo largo de
        la partida y su comportamiento bajo tres protocolos de
        generalización: partidas no vistas, periodos posteriores y mazos
        no observados.
\end{enumerate}

\razon{Los cinco objetivos deben poder verificarse al final del semestre.
Se redactan como acciones con producto observable --- formalizar,
construir, cuantificar, implementar, caracterizar --- y no como intenciones
de desempeño, porque un objetivo del tipo «superar un umbral» quedaría
incumplido si el resultado fuera negativo, cuando un resultado negativo
también responde la pregunta.}

\subsection{Contribución esperada}
\label{subsec:contribucion}

Este trabajo aspira a tres contribuciones. En primer lugar, la
caracterización de una tarea que, hasta donde alcanza esta revisión, no ha
sido abordada en el Pok\'emon Trading Card Game, en un dominio donde la
información oculta y la evolución del conjunto de estrategias dominantes
plantean condiciones distintas a las de los juegos estudiados previamente.
En segundo lugar, un protocolo experimental que trata la fuga de
información como objeto de medición y no únicamente como precaución
declarada. En tercer lugar, un análisis de la calidad de las probabilidades
estimadas, dimensión que Grinsztajn \textit{et al.} identifican como
pendiente en la evaluación de modelos tabulares
\cite{grinsztajn_trees_2022} y que la literatura de predicción de victoria
en juegos ha reportado de forma desigual.

\razon{La contribución se enuncia en términos de aportes verificables y no
de resultados prometidos. Prometer un nivel de desempeño en la Etapa 1
compromete el trabajo con un resultado que aún no se conoce.}

\subsection{Estructura del documento}
\label{subsec:estructura}

La Sección~\ref{sec:relacionados} revisa la literatura en tres ejes
---el problema y su dominio, las técnicas del \textit{track} seleccionado,
y el conjunto de datos--- y cierra con la identificación explícita de las
brechas. La Sección~\ref{sec:baseline} declara el punto de referencia
tomado de la literatura y discute su comparabilidad. La
Sección~\ref{sec:metodologia} presenta la formalización del problema, el
protocolo experimental y el plan de modelos de referencia. La
Sección~\ref{sec:analisistecnico} desarrolla el análisis técnico exigido
por el \textit{track} seleccionado, con énfasis en la prevención de fuga de
información. La Sección~\ref{sec:etica} documenta las consideraciones
éticas preliminares.